\section{Week 7: Reflexive Spaces}

\begin{theorem}[Dual Space (2.10-4)]\label{2.10-4}
   The dual space \( X' \) of a normed space \( X  \) is a Banach Space (whether or not \( X  \) is). 
\end{theorem}

\begin{prop}[ ]
   The dual space of \( \R^{n} \) is \( \R^{n} \).
\end{prop}

\begin{theorem}[3.8-1]\label{Riesz Representation}
    Every bounded linear functional \( f  \) on a Hilbert space \( H  \) can be represented in terms of the inner product , namely, 
    \[  f(x) = \langle x  ,  z  \rangle \]
    where \( z  \) depends on \( f  \), is uniquely determined by \( f  \) and has norm
    \[  \|z \| = \|f \|. \]
\end{theorem}

Recall the definition of Algebraic Reflexivity of a vector space \( X  \) that was introduced in section 2.8.

\begin{definition}[Algebraic Reflexivity]
   A vector space \( X  \) is said to be \textbf{algebraic reflexive} if the canonical mapping \( C : X \to X^{**} \) is surjective. 
\end{definition}

\begin{itemize}
    \item Here we see that \( X^{**} \) is the second algebraic dual space of \( X  \) and the mapping \( C  \) is defined by \( x \mapsto {g}_{x} \) where 
        \[  {g}_{x}(f) = f(x). \ \  \tag{(1) \ \( (f \in X^{*}) \)}; \]
        that is, for any \( x \in X  \) the image is the linear functional \( {g}_{x} \) defined by the above.
    \item If \( X  \) is finite-dimensional, then \( X  \) is algebraic reflexive.
\end{itemize}

Consider a normed space \( X  \) and its dual space \( X' \) (the space of all bounded linear functionals). Observe that \( (X')' \) is the \textbf{second dual space} of \( X  \) (or \textbf{bidual space} of \( X  \)). We define a similar mapping to the above statement, by setting
\[  {g}_{x}(f) = f(x) \ \ \tag{(2), \ \( (f \in X') \)}. \]
The only difference here is that \( f  \) is now bounded and so \( {g}_{x}  \) is bounded. 

\begin{lemma}[Norm of \( g_x \)]
For every fixed \( x  \) in a normed space \( X  \), the functional \( {g}_{x} \) defined above is a bounded linear functional on \( X' \), so that \( {g}_{x} \in X'' \), and has the norm      
\[  \|{g}_{x}\| = \|x\|.  \tag{3}\]
\end{lemma}
\begin{proof}
It is easy to see that \( {g}_{x} \) is linear, and (3) follows from (2) and corollary 4.3-4; that is,
\[  \|{g}_{x}\| = \sup_{\substack{f \in X' \\ f \neq 0}} \frac{ | {g}_{x}(f) |  }{ \|f\| }  = \sup_{\substack{f \in X' \\ f \neq 0 }} \frac{ |f(x)|  }{  \|f \| }  = \|x \|. \tag{4} \]
\end{proof}

\begin{itemize}
    \item This lemma tells us that for every \( x \in X  \), there corresponds a unique bounded linear functional \( {g}_{x} \in X'' \) given by (2) which defines the \textbf{canonical mapping} of \( X  \) into \( X'' \):
        \begin{align*}
            C: X &\longrightarrow X'' \tag{5} \\
                x &\longmapsto {g}_{x}
        \end{align*}
    \item This map is linear and injective and preserves the norm; that is, in terms of mapping the domain to its range, it will be an isomorphism of normed spaces.
\end{itemize}

\begin{lemma}[Canonical Mapping]\label{Canonical Mapping Lemma}
   The canonical mapping \( C  \) given by (5) is an isomorphism of the normed space \( X  \) onto the normed space \( \R(C) \), the range of \( C  \). 
\end{lemma}
\begin{proof}
We can see easily see the linearity of \( C  \). Indeed, we have
\[  {g}_{\alpha x +  \beta y}(f) = f(\alpha x + \beta y) = \alpha f(x)  + \beta f(y) = \alpha {g}_{x}(f) + \beta {g}_{y}(f). \]
In particular, we have \( {g}_{x} - {g}_{y} = {g}_{x -y} \). Thus, we obtain 
\[  \|{g}_{x} - {g}_{y} \| = \|{g}_{x-y}\| = \|x - y\|. \]
This tells us that \( C  \) is an isometry and so it must preserve the norm. Because isometries are injective, \( C  \) is also injective. This can be directly seen from the above formula. Indeed, for any \( x \neq y  \), we have \( {g}_{x} \neq {g}_{y} \) by axiom (N2) in section 2.2. Since we are considering maps from \( X  \) to the range of \( C  \), it follows that \( C  \) is also surjective. Thus, \( C  \) must be bijective.
\end{proof}

\begin{definition}[Embeddable]
    We say that a set \( X  \) is \textbf{embeddable} in a normed space \( Z  \) if \( X  \) is isomorphic with a subspace of \( Z  \).
\end{definition}

\begin{itemize}
    \item From our {\hyperref[Canonical Mapping Lemma]{Canonical Mapping Lemma}}, we can see that \( X  \) is embeddable in \( X'' \) since the range of \( C  \) is a subspace of \( X'' \) and it is isomorphic to \( X  \). 
    \item The mapping \( C  \) is called the \textbf{canonical embedding} in this case.
    \item In general, the range of \( C  \) will be a proper subspace of \( X'' \). 
    \item If the range of \( C  \) happens to be a surjection, then we will call it a different name which will be defined below.
\end{itemize}

\begin{definition}[Reflexivity]
    A normed space \( X  \) is said to be \textbf{reflexive} if 
    \[  \mathcal{R}(C) = X'' \]
    where \( C : X \to X'' \) is the canonical mapping given by (5) and (2).
\end{definition}

\begin{itemize}
    \item If \( X  \) is reflexive, it is isomorphic (and hence isometric) with \( X'' \) by {\hyperref[Canonical Mapping Lemma]{Lemma 4.6-2}}.  
    \item However, the converse does not hold.
    \item Completeness does not necessarily imply reflexivity.
    \item However, the converse does hold which will be outlined below.
\end{itemize}

\begin{theorem}[Completeness]
    If a normed space \( X  \) is reflexive, it is complete (hence a Banach space).
\end{theorem}
\begin{proof}
    Since \( X'' \) is the dual space of \( X' \), it is complete by {\hyperref[2.10-4]{Theorem 2.10-4}}. Reflexivity of \( X  \) means that \( \mathcal{R}(C) = X''\). Completeness of \( X  \) now follows from that of \( X'' \) by {\hyperref[Canonical Mapping Lemma]{Lemma 4.6-2}}.  
\end{proof}

\begin{itemize}
    \item An example of a reflexive space is \( \R^{n} \) which follows from example 2.10-5.
    \item For finite dimensions, every linear functional on a general normed space \( X  \) is bounded and so \( X' = X^{*} \) which implies further that 
\end{itemize}

\begin{theorem}[Finite Dimension]
    Every finite-dimensional normed space is reflexive.
\end{theorem}

\begin{itemize}
    \item \( l^{p} \) with \( 1 < p < + \infty  \) is reflexive.
    \item \( L^{p}[a,b]  \) with \( 1 < p < + \infty  \) is reflexive.
    \item Non-examples are \( C[a,b] \), \( l^{1} \), \( L^{1}[a,b] \), \( l^{\infty } \) and the subspaces \( c  \) and \( {c}_{0}  \) of \( l^{\infty } \) where \( c  \) is the space of all convergent sequence of scalars and \( {c}_{0} \) is the space of all sequences of scalars converging to zero.
\end{itemize}

\begin{theorem}[Hilbert Space]
   Every Hilbert Space \( H  \) is reflexive. 
\end{theorem}

\begin{proof}
Our goal is to show that the canonical mapping \( C : H \to H'' \) is surjective by showing that for every \( g \in H'' \), there exists an \( x \in H  \) such that \( g = C x  \). Define \( A : H' \to H  \) by \( Af = z  \), where \( z  \) is given the {\hyperref[Riesz Representation]{Riesz Representation}} \( f(x) = \langle x , z  \rangle  \). From this, we see that \( A  \) is bijective and isometric. Note that \( A  \) is conjugate linear, as we see from (16) from section 4.5. Now, \( H' \) is complete by {\hyperref[2.10-4]{2.10-4}} and a Hilbert space with inner product defined by
\[  \langle {f}_{1}  ,  {f}_{2} \rangle_1 = \langle A {f}_{2} ,  A {f}_{1} \rangle. \]
It can be easily shown that \( \langle  ,  \rangle_1  \) is an inner product on \( H  \).

Let \( g \in H'' \) be arbitrary. Let its Riesz Representation be 
\[  g(f) = \langle f , {f}_{0} \rangle_1 = \langle A {f}_{0} ,  A f \rangle. \]
We now remember that \( f(x) = \langle x  , z  \rangle  \) where \( z = A f  \). Writing \( A {f}_{0} = x  \), we have 
\[  \langle A {f}_{0}  ,  A f  \rangle = \langle x  , z  \rangle = f(x). \]
Hence, we have \( g(f) = f(x) \) and so \( g = C x  \) by definition of \( C  \). Hence, \( C  \) is surjective and so \( H  \) must be reflexive.
\end{proof}

\begin{itemize}
    \item Separability of the dual of a normed space implies separability of \( X  \) (note the converse is not generally true).
    \item  If a normed space \( X  \) is reflexive, then \( X'' \) is isomorphic with \( X  \) by {\hyperref[4.6-2]{4.6-2}}, and so \( X  \) implies separability of \( X'' \). Thus, the dual space \( X' \) must also be separable.
    \item \textit{A separable normed space \( X  \) with a nonseparable dual space \( X' \) cannot be reflexive.} 
\end{itemize}

\begin{eg}[\( l^1  \) is not reflexive]
    Note that \( l^1  \) is separable by 1.3-10, but the dual of \( l^1  \) which is \( l^{\infty } \) is not via 2.10-6 and 1.3-9.
\end{eg}
